% TOP_PROBLEM: {"id": 30006835, "title": "Frey-Mazur Conjecture", "category_id": 1, "category": {"id": 1, "name": "number_theory", "display_name": "Number Theory", "description": "Properties of integers, prime numbers, Diophantine equations.", "slug": "number-theory", "order_index": 1, "created_at": "2026-07-31T15:26:25.670Z"}, "status": "open"}
% FIRST_POSTED: 2026-09-27
% !TeX program = lualatex
% ATTEMPT_STATUS: unresolved
% Research continuation by GPT_6_Astra_Ultra, 27 September 2026.
\documentclass[11pt]{article}
\usepackage[margin=25mm]{geometry}
\usepackage{fontspec}
\setmainfont{Latin Modern Roman}
\usepackage{amsmath,amssymb,mathtools}
\usepackage{unicode-math}
\setmathfont{Latin Modern Math}
\usepackage{xurl,hyperref}
\hypersetup{colorlinks=true,urlcolor=blue}
\setcounter{secnumdepth}{0}
\setlength{\parindent}{0pt}
\setlength{\parskip}{5pt}
\setlength{\emergencystretch}{3em}
\begin{document}
\section*{\texorpdfstring{Frey-Mazur Conjecture}{Frey-Mazur Conjecture}}\label{30006835}
\subsection{Definitions and mathematical statement}
For an elliptic curve over ${\mathbb{Q}}$, $E[p]$ denotes the two-dimensional ${\mathbb{F}}_p$-space of geometric $p$-torsion, with its action of $G_{{\mathbb{Q}}}=\operatorname{Gal}(\overline{{\mathbb{Q}}}/{\mathbb{Q}})$. A ${\mathbb{Q}}$-isogeny is a nonzero morphism of elliptic curves over ${\mathbb{Q}}$ preserving the origin, necessarily with finite kernel. The selected Frey--Mazur assertion is
\[
 p>17,\qquad E[p]\cong E'[p]\text{ as }G_{{\mathbb{Q}}}\text{-modules}
 \quad\Longrightarrow\quad E\sim_{{\mathbb{Q}}}E'.
\]
The isomorphism is not additionally required to preserve the Weil pairing. The explicit cutoff $17$ is part of the target; existence of some unspecified larger universal cutoff is weaker.
\par\smallskip\textit{Formulation and concept references:} \href{https://annals.math.princeton.edu/2016/184-3/p02}{[S1]}.

\subsection{Short English statement}
For every prime greater than seventeen, does identical Galois-module structure on two rational elliptic curves' p-torsion force the curves to be isogenous over the rationals?
\subsection{Research attempt: trace rigidity and the missing uniform bound}
\textit{GPT-6 Astra Ultra, 27 September 2026. Status: UNRESOLVED.}

An isomorphism of residual Galois modules imposes infinitely many Frobenius congruences. Small traces turn some of these congruences into equalities. I make that argument exact, prove a fixed-pair finiteness consequence, and construct nonisogenous twists agreeing at any prescribed finite set of good primes. The construction is a countertest to finite-trace shortcuts, not a counterexample to the residual-module assertion.

\paragraph{A rigorous small-prime consequence.}
At a prime $\ell\ne p$ of good reduction for both curves, put
\[
 a_\ell(E)=\ell+1-\#E(\mathbb F_\ell).
\]
The characteristic polynomial of Frobenius on $E[p]$ is the reduction modulo $p$ of $T^2-a_\ell(E)T+\ell$. Therefore the hypothesized module isomorphism implies
\[
 a_\ell(E)\equiv a_\ell(E')\pmod p.                        \tag{1}
\]
The determinant congruence contains no additional distinction at these primes: both determinants are $\ell$ modulo $p$. No preservation of the Weil pairing is needed for (1), in agreement with the formulation.

Hasse's estimate gives $|a_\ell(E)|,|a_\ell(E')|\le2\sqrt\ell$, so
\[
 |a_\ell(E)-a_\ell(E')|\le4\sqrt\ell.
\]
If $p>4\sqrt\ell$, the only multiple of $p$ in this interval is zero. Consequently
\[
 a_\ell(E)=a_\ell(E')\quad
 \text{for every common good }\ell\ne p\text{ with }\ell<p^2/16.
                                                               \tag{2}
\]
The strict inequality is intentional. At an endpoint one cannot conclude that a multiple of $p$ with absolute value at most $p$ is zero. Also, the restriction to common good primes cannot be omitted, since the displayed quadratic Frobenius polynomial was derived under good reduction.

\paragraph{For a fixed pair, only finitely many exceptional primes.}
Faltings' isogeny theorem, together with the usual Frobenius-density comparison, implies that two elliptic curves over $\mathbb Q$ with equal traces at all common good primes are $\mathbb Q$-isogenous [E1]. Therefore a fixed nonisogenous pair has a common good prime $\ell_0$ with
\[
 d=a_{\ell_0}(E)-a_{\ell_0}(E')\ne0.
\]
If its $p$-torsion modules are isomorphic and $p\ne\ell_0$, then $p\mid d$ by (1). There are only finitely many such primes, and each satisfies
\[
 p\le |d|\le4\sqrt{\ell_0}.                               \tag{3}
\]
The possible prime $p=\ell_0$ is simply retained as one extra exception because (1) was not used there. This proves finiteness for each fixed pair without proving a uniform cutoff.

The distinction is the location of the quantifier. In (3), the witness $\ell_0$ depends on $E,E'$. Its size can grow as the pair varies. No step gives $\ell_0\le17^2/16$, or any other uniform bound sufficient to replace (3) by the constant $17$. A finite computation for each curve in a bounded database likewise leaves curves of larger height or conductor outside its scope.

\paragraph{Finite agreement cannot identify the isogeny class.}
Here is an explicit construction. Start with a non-CM elliptic curve $E/\mathbb Q$, and let $S$ be any finite set of odd good primes. Choose a positive integer
\[
 M=8\prod_{\ell\in S}\ell,
 \qquad D=M^2+M+1.
\]
Then $M^2<D<(M+1)^2$, so $D$ is not a rational square, and $D\equiv1\pmod\ell$ for every $\ell\in S$. The quadratic twist $E^{(D)}$ has good reduction at those primes and
\[
 a_\ell(E^{(D)})=\left(\frac D\ell\right)a_\ell(E)
               =a_\ell(E)\qquad(\ell\in S).               \tag{4}
\]
It is not necessary for $D$ to be squarefree: replacing it by its squarefree part defines the same quadratic twist, and the square classes modulo primes in $S$ remain unchanged.

To justify nonisogeny, let $\phi:E_{\overline{\mathbb Q}}\to E^{(D)}_{\overline{\mathbb Q}}$ be the standard twist isomorphism. Because $E$ is non-CM, $\operatorname{End}_{\overline{\mathbb Q}}(E)=\mathbb Z$, so every geometric homomorphism from $E$ to $E^{(D)}$ is an integral multiple of $\phi$. A Galois element nontrivial on $\sqrt D$ sends $\phi$ to $-\phi$. Hence no nonzero multiple is Galois invariant. There is no nonzero $\mathbb Q$-homomorphism and thus no $\mathbb Q$-isogeny.

For a concrete input one may take $E:y^2=x^3-x+1$. Its discriminant is $-16\cdot23$ and its $j$-invariant is $-6912/23$. A CM $j$-invariant is an algebraic integer; this nonintegral rational value certifies that the curve is non-CM [E2]. Thus (4) gives actual nonisogenous pairs with any prescribed finite collection of matching good traces, not just a hypothetical construction.

These twists are \emph{not} claimed to have isomorphic $p$-torsion for any chosen large $p$. They show that using only the finite equalities furnished by (2), without exploiting the rest of the residual representation, cannot establish the desired global conclusion. For any fixed finite test list there are nonisogenous pairs passing that list.

\paragraph{A residual self-twist has stronger constraints.}
The twist relation is
\[
 \overline\rho_{E^{(D)},p}\cong
 \overline\rho_{E,p}\otimes\chi_D.
\]
If the two residual modules were isomorphic, then at every common good $\ell\ne p$ with $\chi_D(\ell)=-1$ we would have
\[
 a_\ell(E)\equiv-a_\ell(E)\pmod p,
 \qquad 2a_\ell(E)\equiv0\pmod p.
\]
For odd $p$, a nonzero such trace with $p>2|a_\ell(E)|$ disproves the isomorphism. This is an effective obstruction for a specified pair once a suitable prime is found. It also shows why the construction (4) is consistent with Frey--Mazur: it controls the split primes in a finite set but says nothing about the inert primes that can expose the twist.

An argument based only on congruence of characteristic polynomials has another limitation. Equality of traces generally identifies semisimplifications of representations under suitable density hypotheses, not every possible extension class. The actual hypothesis of the problem is an isomorphism of modules and is stronger. The partial implications above all run from that hypothesis to traces; no converse from a finite trace list has been used.

\paragraph{What a sufficient uniform trace theorem would say.}
For this route, a sufficient additional statement would be: whenever $p>17$ and two nonisogenous rational elliptic curves have isomorphic $p$-torsion, there exists a common good prime $\ell\ne p$ with $a_\ell(E)\ne a_\ell(E')$ and $4\sqrt\ell<p$. Equations (1)--(2) would contradict it immediately. This statement is formulated conditionally on residual isomorphism; the twist construction shows why one cannot replace it by an assertion about every nonisogenous pair without the residual condition.

No such uniform theorem is proved here. Standard effective comparisons whose bounds depend on conductors or heights do not automatically have this strength: those quantities can be unbounded while $p$ is held fixed. Nor would establishing an unspecified universal bound $C$ settle the printed cutoff $17$ without analyzing the primes between $17$ and $C$.

The inherited source [S1] proves a geometric function-field analogue for families over complex curves, with a bound depending on the base curve's gonality and with the stated nonisotriviality assumptions. It does not prove the rational-number-field assertion with cutoff $17$. Transferring its conclusion would require an additional arithmetic argument, not a change of notation from a function field to $\mathbb Q$.

\paragraph{Verification and outcome.}
The local exact checks count points on the displayed curve and its chosen twist over a finite list of primes, verify (4), and enumerate the integer bounds underlying (2). The proof of nonisogeny is the Galois action on geometric homomorphisms; finite point counts alone are not used to certify it. The completed output is small-prime trace rigidity, fixed-pair finiteness, and an explicit limitation on finite-trace tests. The first unresolved step is uniform arithmetic rigidity strong enough to force isogeny for every $p>17$. The conjecture remains unresolved by this attempt.

\subsection{Sources}
\begin{itemize}
\item[S1]Benjamin Bakker and Jacob Tsimerman, "p-torsion monodromy representations of elliptic curves over geometric function fields," Annals of Mathematics 184 (2016), 709-744.
\url{https://annals.math.princeton.edu/2016/184-3/p02}.
\textit{Formulation source.}
\item[E1]G. Faltings, \textit{Endlichkeitss\"atze f\"ur abelsche Variet\"aten \"uber Zahlk\"orpern}, Inventiones Mathematicae 73 (1983), 349--366. Used for the isogeny comparison after equality of all good Frobenius traces, not a uniform residual bound.
\url{https://doi.org/10.1007/BF01388432}.
\item[E2]J. H. Silverman, \textit{The Arithmetic of Elliptic Curves}, second edition, Springer, 2009. Standard Frobenius, twisting, Hasse-bound and CM integrality facts used above.
\url{https://doi.org/10.1007/978-0-387-09494-6}.

\end{itemize}
\end{document}
